% ##############################################################################
  \section{Runtime Composition: Making the UI Tell the Truth}
  This section describes how the web tier turns the UI contract into pixels
  without giving the browser a second job as an authorization engine.

% ==============================================================================
  \subsection{Abilities: one source of truth, exported for UX}
  Authorization is enforced server-side, but the UI still needs to know
  what to show. We use CASL \cite{casl_docs} to represent abilities and explicitly
  avoid duplicating rules in the frontend.

  \paragraph{Design.}
  \begin{itemize}
    \item The server builds the canonical ability from session-backed profile
      data (modules + role).

    \item The server serializes the ability rules using CASL's packed format.

    \item The client unpacks these rules to construct an identical ability
      object for UI/UX gating only.
  \end{itemize}

  \begin{samepage}
  \begin{lstlisting}[style=typescript,caption={Server-authoritative abilities: rules are built server-side and packed for client use. The backend remains the enforcement point.},label={lst:casl_pack}]
import { packRules } from '@casl/ability/extra'

// GET /api/auth/user/ability
const ability = await getUserAbility()
return { rules: packRules(ability.rules) }
\end{lstlisting}
  \end{samepage}

  \paragraph{Two-layer protection (by design).}
  The client uses the ability to hide inaccessible UI elements and prevent
  dead ends. Server components and API routes still enforce
  authorization. In other words: the UI is truthful, but it is not trusted.

% ==============================================================================
  \subsection{From modules to navigation: composition, not conditionals}
  The module list in \texttt{view\_as.modules} drives the UI skeleton:

  \begin{itemize}
    \item The sidebar renders only modules in the list.

    \item Module pages are protected server-side (redirects/404) when the
      user lacks \texttt{view} permission on that subject.

    \item Cross-module UI behaviors (top bar, settings entry points,
      feedback widgets) are controlled by \texttt{common} layout
      fragments.
  \end{itemize}

  This yields an important property: tenant onboarding becomes a
  configuration change, not a UI rewrite. The product behaves like a
  platform instead of a branching forest of \texttt{if (tenant == \ldots)}.

% ==============================================================================
  \subsection{Lazy module layout fetching and deep merging}
  Shipping an entire tenant UI contract on every request is wasteful.
  Instead, we lazy-load module-specific configuration on demand.

  \paragraph{Mechanism.}
  The web tier exposes:
  \begin{center}
    \texttt{/api/auth/user/module-layout?module=<name>}
  \end{center}
  This endpoint:
  \begin{enumerate}
    \item constructs a fallback default layout for \texttt{common} and (optionally)
      the requested module,

    \item checks module access via server-side CASL for non-\texttt{common}
      modules,

    \item merges the tenant-provided layout fragments from \texttt{/v2/me}
      over defaults using a recursive deep-merge.
  \end{enumerate}

  \begin{samepage}
  \begin{lstlisting}[style=typescript,caption={Merge semantics: nested objects recursively merge, primitives/arrays replace. This enables partial overrides without forcing tenants to copy entire configs.},label={lst:merge_semantics}]
function deepMerge(target, source) {
  const result = { ...target }
  for (const key in source) {
    const s = source[key]
    const t = result[key]
    if (s !== null && s !== undefined) {
      if (isObject(s) && isObject(t)) result[key] = deepMerge(t, s)
      else result[key] = s
    }
  }
  return result
}
\end{lstlisting}
  \end{samepage}

  \paragraph{Why deep-merge matters.}
  Deep-merge keeps tenant config surgical. A tenant can override \texttt{common.sideBar.defaultCollapsed}
  without duplicating the rest of the UI definition, keeping
  configuration maintainable and reviewable.

% ==============================================================================
  \subsection{Truthful UX as an invariant}
  We treat ``no broken doors'' as an invariant:

  \begin{itemize}
    \item If the backend denies access to a module, the module does not
      appear in navigation.

    \item If a user tries a direct URL, server-side checks redirect/deny.

    \item If config asks for a module the user lacks, CASL checks
      prevent module config from being served.
  \end{itemize}

  Put differently: the UI is not permitted to hallucinate. (That is
  reserved for marketing.)
